\documentclass{article}
\usepackage[T1]{fontenc}
\usepackage[utf8]{inputenc}
\usepackage{multicol}
\begin{document}

\begin{multicols}{2}

\section{Left Hand Column}
The left column begins with this section. It carries the distinctive marker
\texttt{columnsentinel-left}, which the tests use to confirm that text placed
in the left column of a two-column layout is extracted. Additional filler
prose about botany and seasonal rainfall follows so that the left column
achieves a reasonable height and the layout genuinely flows across the two
columns rather than collapsing into a single narrow strip.

\section{Right Hand Column}
The right column carries its own distinctive marker
\texttt{columnsentinel-right}. This text appears on the same page but in the
second column, which is the point of this fixture: to verify that two-column
layout is still extracted as readable text, with content from both columns
present in the output even though their physical positions on the page differ.

\end{multicols}

\end{document}
